\section{The exact minimum budget for a counterexample}
\label{sec:exact}
Fix any prime $p$, depth $\rho=1$, and tolerance $\epsilon\ge0$. A family
\emph{blocks} the first $H$ shells if every allowed root of unity in those shells
has total positive excess strictly greater than $\epsilon$. Write
$C_p(\epsilon,H)$ for the least witness budget of such a family, allowing
arbitrarily many integral polynomials of arbitrary degrees over
$\overline{\mathbb Q}_p$. The minimum is attained. The following formula
computes it without searching over polynomials or coefficient fields.

\begin{theorem}\label{th:cost}
With $e_j=\varphi(p^j)$ define
\[
 U_h(m)=
 \begin{cases}
 \displaystyle\frac{1}{e_h}\sum_{j=1}^{h-1}e_j
                  \left\lfloor\frac{m}{e_j}\right\rfloor,&0\le m<e_h,\\[5pt]
 +\infty,& m\ge e_h .
 \end{cases}
\]
Let $m_h$ be the least nonnegative integer with $U_h(m_h)>\epsilon$. Then
\[
 \boxed{\displaystyle C_p(\epsilon,H)=\sum_{h=1}^H\frac{m_h}{e_h}.}
\]
Thus $B<C_p(\epsilon,H)$ guarantees success somewhere in shells $1$ through $H$.
At equality an explicit counterexample already exists with every witness at zero.
The statement also holds for a single polynomial if the excess in shell $h$ is
measured over its moving Gauss valuation $G_{r_h}$ and its budget is $G_1$.
\end{theorem}

\subsection*{A single shell}
Set $z_a=\xi^a-1$ for a fixed primitive $p^h$-th root $\xi$ and let $a$ run through
the units modulo $p^h$. Given $m$ algebraic roots $\alpha$ of valuation $r_h$,
counted with multiplicity, consider
\[
 \min_{a}\sum_\alpha \big(v(z_a-\alpha)-r_h\big).
\]
Its largest possible value is $U_h(m)$. To see the upper bound, replace each
$\alpha$ by a $z_b$ maximizing $v(z_b-\alpha)$ on this finite set.
For every $a$ the ultrametric inequality gives
$v(z_a-z_b)\ge\min\{v(z_a-\alpha),v(z_b-\alpha)\}=v(z_a-\alpha)$.
This replacement can only raise the minimum. We may therefore reason about a
multiset of $m$ unit residues $b$.

For a given $a$ write $N_j(a)$ for the number of these residues congruent to $a$
modulo $p^j$. If $a$ itself is unoccupied, the order of $\xi^{a-b}$ gives
\[
 v(z_a-z_b)=r_h p^{v_p(a-b)},\qquad
 \sum_b(v(z_a-z_b)-r_h)=r_h\sum_{j=1}^{h-1} e_j N_j(a).
\]
The second equality expands $p^t-1=\sum_{j=1}^t(p-1)p^{j-1}$.
Choose a least occupied unit class modulo $p$, then a least occupied child
modulo $p^2$, and continue. The first choice has at most $\lfloor m/(p-1)\rfloor$
occupants, and at each subsequent step one divides by $p$ and rounds down.
Consequently $N_j(a)\le\lfloor m/e_j\rfloor$ at every depth. If $m<e_h$, the
terminal residue is unoccupied. This proves the upper bound simultaneously at
all depths.

For the reverse bound, place the roots at the first $m$ positive unit exponents
in increasing order, omitting multiples of $p$. In every block of $e_j$
successive unit exponents, each unit class modulo $p^j$ occurs once. Thus every
$N_j(a)\ge\lfloor m/e_j\rfloor$. Occupied leaves give infinite valuation and
the other leaves give at least the required sum. For $m\ge e_h$, putting a root
at every leaf makes the minimum infinite. This includes the shell of order two
and the case $m=0$ by the empty-sum convention. The one-shell argument is a
cyclotomic instance of the fixed-divisor extremal principle underlying
Bhargava's $P$-orderings \cite{bhargava}. The global budget is a separate constraint.

\subsection*{Passing from one shell to every shell}
For a general family, take all roots of all the $g_i$, with multiplicities, and
let $M_h$ count the ones with valuation $r_h$. Factorization, or
multiplicativity of the Gauss valuation, yields
\[
 \sum_{h=1}^H M_h r_h\ \le\ \sum_i\big(G_1(g_i)-G_0(g_i)\big)
 \ \le\ \sum_i v(g_i(x_i)).
\]
In fact a root of valuation $w$ contributes $\min(1,w)-\min(0,w)$ to the
middle expression, which also handles roots of negative valuation, and
all $G_0(g_i)$ are nonnegative. For a primitive point in shell $h$ we have
$t_i\ge G_1(g_i)\ge G_{r_h}(g_i)$. The total positive excess is bounded by
$\sum_i (v(g_i(z_a))-G_{r_h}(g_i))$. In that sum every root with valuation
different from $r_h$ contributes zero, by the strict ultrametric equality.
The single-shell upper bound shows that blocking requires $M_h\ge m_h$.
Summing the charges proves the lower bound on $C_p$.

For equality, use one monic polynomial per shell, with roots at the first
$m_h$ ordered unit exponents as above. Its value at zero has valuation exactly
$m_hr_h$. On its designated shell its excess is at least $U_h(m_h)>\epsilon$.
The positive part prevents any saving from the other polynomials from canceling
this excess. All coefficients lie in $\mathbb Q_p(\zeta_{p^H})$. Multiplying
the polynomials gives the corresponding attainment for the moving Gauss
valuation, because its excess is additive and each summand is nonnegative.
\hfill$\square$

\paragraph{An exact contrasting case.}
For coefficients in a finite unramified extension of $\mathbb Q_p$,
the bound in the preceding section guarantees zero excess within the first
$\lfloor B\rfloor+1$ shells. It is optimal for every finite tolerance.
Indeed $\Phi_{p^h}(1+T)$ vanishes throughout shell $h$ and has value
$p$ at zero. Taking one such polynomial per shell costs exactly $H$ for $H$
shells. For unrestricted algebraic coefficients, subdividing the cost within
one shell is precisely what the floor sums measure.
